\paragraph{Audio and video.} The same pattern holds in the other modalities. Audio attacks were judged, like the image attacks, at the pair-level $10^{-6}$ threshold of each method, taken from the first audio retrieval run; those thresholds differ from the re-run reported in Section~\ref{sec:fp-res-audio} by at most \val{fp:sec/audio/thdelta}, and re-scoring every CLAP attack at the re-run thresholds changes \val{fp:sec/audio/rescore_changed} outcomes. Video attacks were judged at a different threshold, described below. For audio, projected gradient descent on the CLAP embedding of \val{fp:sec/audio/n} registered tracks evaded CLAP for every track at every budget we tried, down to a target signal-to-noise ratio of 40\,dB (achieved median \val{fp:sec/audio/snr/40}\,dB). The perturbation did not transfer to the recording-level fingerprints: Chromaprint matched every attacked track at all budgets, and the ISCC Audio-Code and audfprint lost \val{fp:sec/audio/pgd/20/ISCC-Audio-64/pct}\% and \val{fp:sec/audio/pgd/20/audfprint/pct}\% only at the loudest budget (20\,dB target), where the added signal acts as noise rather than as a targeted attack (\val{fp:sec/audio/pgd/40/ISCC-Audio-64/pct}\% and \val{fp:sec/audio/pgd/40/audfprint/pct}\% at 40\,dB). An attacker does not need gradients against those systems, however: a search over pitch and tempo changes found the smallest change that broke each match, and the median was \val{fp:sec/audio/mag/Chromaprint/pitch} semitones or a tempo factor of \val{fp:sec/audio/mag/Chromaprint/tempo} for Chromaprint and \val{fp:sec/audio/mag/audfprint/pitch} semitones or \val{fp:sec/audio/mag/audfprint/tempo} for audfprint and the ISCC Audio-Code, where 0.25 semitones and 1.02 are the smallest steps of the grid. Only CLAP needed audible changes (\val{fp:sec/audio/mag/CLAP/pitch} semitones, tempo \val{fp:sec/audio/mag/CLAP/tempo}). Every audio fingerprint in this study is therefore evaded by a small change, either by gradients (CLAP) or by a pitch shift (all others).

For video, frame-wise projected gradient descent against DINOv2-S and SSCD on \val{fp:sec/video/n} registered clips evaded both frame-level neural descriptors, in sequence and mean-pooled form, for every clip at 4 grey levels, and transferred to none of the hashes: vPDQ, TMK+PDQF and both ISCC Video-Codes matched every attacked clip at 4 and 8 grey levels (judged at each method's query-level 1\% threshold). The video calibration set does not resolve a pair-level $10^{-6}$ target, so every video attack used that threshold, which let through \val{fp:video/fprq/min}\% to \val{fp:video/fprq/max}\% of held-out negative queries (Section~\ref{sec:fp-res-video}). The video results therefore do not establish attack performance at a pair-level false-match rate. Each hash's pair-level $10^{-4}$ threshold is lower than its query-level threshold, so the hashes' matches hold at that target too; the per-clip scores of the neural descriptors were not kept, so their evasion was not re-scored at it. Simple edits again did what gradients could not. A centred crop keeping 90\% of the width and height, the first step of the grid, broke the median match for vPDQ, TMK+PDQF, both ISCC Video-Codes and the sequence-matched DINOv2-S; only SSCD and mean-pooled DINOv2-S held at 90\% and broke at a median of \val{fp:sec/video/magnitude/SSCD-mixup-seq/crop_keep_median}. Playback speed-up was the reverse: TMK+PDQF and the 256-bit ISCC Video-Code still matched all \val{fp:sec/video/magnitude/TMK+PDQF/speed_never} clips at twice the speed (the 64-bit code \val{fp:sec/video/magnitude/ISCC-Video-64/speed_never}), because dropping frames leaves the remaining frames' hashes intact, while the sequence-matched neural descriptors broke at a median factor of \val{fp:sec/video/magnitude/DINOv2-S-seq/speed_median}. No video fingerprint survived both a small crop and a white-box attack.
